\section{Quality Metrics for Point Cloud Compression}
\label{sec:metrics}

\subsection{Point-to-Point PSNR (D1)}
\label{ssec:d1-psnr}

The point-to-point PSNR (D1) is the standard metric used in MPEG point cloud
codec evaluations~\cite{tian2017geometric}.
For a reconstructed cloud $\hat{\pc}$ and original $\pc$, two directional
mean squared errors are computed:
\begin{align}
  e_{\text{fwd}} &= \frac{1}{|\pc|}\sum_{p \in \pc} \min_{\hat{p} \in \hat{\pc}} \|p - \hat{p}\|^2, \\
  e_{\text{rev}} &= \frac{1}{|\hat{\pc}|}\sum_{\hat{p} \in \hat{\pc}} \min_{p \in \pc} \|\hat{p} - p\|^2.
\end{align}
The D1 PSNR is:
\begin{equation}
  \PSNR_{\text{D1}} = 10\log_{10}\!\left(\frac{d_{\max}^2}{\max(e_{\text{fwd}},\, e_{\text{rev}})}\right),
\end{equation}
where $d_{\max}$ is the diagonal of the axis-aligned bounding box of the
original cloud, used as the peak signal value.
Taking the maximum of the two directions ensures that both missing coverage
(forward) and spurious points (reverse) are penalised.

\paragraph{Limitation of D1.}
D1 averages over all $N$ points with equal weight, regardless of their
geometric role.
A point cloud in which 90\% of points lie on flat surfaces and 10\% on edges
will have its D1 PSNR dominated by flat-region fidelity.
A codec that faithfully reconstructs the flat regions while systematically
erasing all edges can report high D1 PSNR that does not reflect the
perceptual quality loss.

This limitation motivates the curvature-stratified metric introduced in
Section~\ref{ssec:curvature-stratified}.

\subsection{Point-to-Plane PSNR (D2)}

Point-to-plane PSNR (D2) replaces the Euclidean nearest-neighbour distance
with its projection onto the reference surface normal $\mathbf{n}_p$:
\begin{equation}
  e_{\text{D2}} = \frac{1}{|\pc|}\sum_{p \in \pc}
    \left\langle p - \hat{p}^*,\, \mathbf{n}_p \right\rangle^2,
  \qquad \hat{p}^* = \arg\min_{\hat{p}} \|p - \hat{p}\|^2.
\end{equation}
This is perceptually motivated: displacements tangential to the surface are
less visible than those in the normal direction.
D2 requires reliable normal estimation at each reference point; for voxelized
clouds at integer positions, normals are estimated from the same
$k$-NN covariance analysis as curvature.
D2 correlates with D1 for most evaluation scenarios~\cite{tian2017geometric}; this work uses D1
as the primary metric for reproducibility.

\subsection{Curvature-Stratified PSNR (This Work)}
\label{ssec:curvature-stratified}

The curvature-stratified PSNR metric introduced in this thesis addresses the
averaging limitation of D1 by evaluating quality separately within surface
strata defined by local curvature.
The metric is applied throughout Chapters~\ref{ch:oversmoothing}
and~\ref{ch:results}.

\paragraph{Step 1: Curvature estimation.}
For each point $p$ in the \emph{original} cloud $\pc$, compute the surface
variation $\kappa(p)$ using the $k = 30$ nearest neighbours
(Section~\ref{ssec:curvature-theory}).

\paragraph{Step 2: Percentile stratification.}
Sort all $|\pc|$ points by their curvature and compute the 33rd and 67th
percentiles $\kappa_{33}$ and $\kappa_{67}$.
Partition into three strata:
\begin{align}
  \mathcal{S}_{\text{flat}} &= \{p \in \pc : \kappa(p) \leq \kappa_{33}\}, \\
  \mathcal{S}_{\text{mid}}  &= \{p \in \pc : \kappa_{33} < \kappa(p) \leq \kappa_{67}\}, \\
  \mathcal{S}_{\text{edge}} &= \{p \in \pc : \kappa(p) > \kappa_{67}\}.
\end{align}
Each stratum contains exactly $|\pc|/3$ points (up to rounding).
Percentile-based thresholds are adaptive: they are specific to each individual
cloud, ensuring that the ``edge'' stratum always identifies the most curved
third of \emph{this} cloud, regardless of its overall curvature level.
This avoids the dataset dependence of a fixed absolute threshold.

\paragraph{Step 3: Per-stratum reverse PSNR.}
For each stratum $\mathcal{S} \in \{\mathcal{S}_{\text{flat}}, \mathcal{S}_{\text{mid}}, \mathcal{S}_{\text{edge}}\}$:
\begin{equation}
  \PSNR_{\mathcal{S}} = 10\log_{10}\!\left(
    \frac{d_{\max}^2}{\frac{1}{|\mathcal{S}|}\sum_{p \in \mathcal{S}}
      \min_{\hat{p} \in \hat{\pc}} \|p - \hat{p}\|^2}
  \right).
  \label{eq:stratified-psnr}
\end{equation}

\paragraph{Why the reverse direction.}
The curvature label is computed on the \emph{original} cloud, where curvature
is geometrically meaningful and stable.
In the \emph{reverse} direction (Equation~\ref{eq:stratified-psnr}),
the outer sum runs over original points and the inner minimisation is over
reconstructed points.
This measures: \textit{for each original edge point, how well is it covered
by the reconstruction?}

The forward direction (reconstructed-to-original) would require curvature
labels on the decoded cloud, where oversmoothing has already reduced curvature
at edge regions.
A decoded cloud that completely erases an edge will have \emph{no} decoded
points in the edge region, so the forward direction would simply skip those
points, producing optimistic edge PSNR.
The reverse direction is immune to this: it always queries the coverage of
original edge points, regardless of what the decoder produced.

\paragraph{Degradation gap.}
The primary diagnostic is:
\begin{equation}
  \Delta_{\text{smooth}} = \PSNR_{\text{flat}} - \PSNR_{\text{edge}} \geq 0,
\end{equation}
which is positive when flat regions are reconstructed better than edges
(the oversmoothing signature).
A post-processing correction that reduces $\Delta_{\text{smooth}}$ is
specifically targeting the most harmful codec artifact.

\subsection{Bjontegaard Delta (BD) Metrics}
\label{ssec:bd}

Bjontegaard Delta metrics~\cite{bjontegaard2001calculation} summarise codec
comparisons over an entire rate-distortion curve, rather than at a single
operating point.
Given two R-D curves (e.g., baseline PCGCv2 and post-processed),
BD-PSNR is the average PSNR gain at the same bitrate, and BD-rate is the
average bitrate saving at the same PSNR, each computed by integrating the
difference of piecewise cubic polynomial fits to the curves.

In this work, BD-PSNR is computed over the three evaluated rate points
(r2, r4, r6), providing a single-number summary of the post-processing
improvement across the evaluated rate range.
